% Vietnamese translation for OI Public Library
% Source-File: IOI中国国家候选队论文/2004/杨思雨.ppt
\documentclass[11pt]{article}
\usepackage{oipl-vn}

\newcommand{\Oh}{\mathcal{O}}

\title{Các thao tác cơ bản và ứng dụng của cây splay\\{\large Ghi chú theo slide}}
\author{Yang Siyu}
\date{IOI 2004}

\begin{document}
\maketitle

\origtitle{Basic operations and applications of splay trees}
\sourcepath{IOI 2004 / Yang Siyu.ppt}

\section{Từ cây tìm kiếm nhị phân đến splay}

Các slide đối chiếu cây tìm kiếm nhị phân thông thường, có thể suy biến tới
\(\Oh(n)\), với cây splay có chi phí khấu hao \(\Oh(\log n)\). Thao tác chủ
đạo là đưa nút vừa truy cập lên gốc bằng các phép quay.

\section{Ba kiểu quay}

Bài trình bày liệt kê ba trường hợp chuẩn:
\begin{itemize}
  \item Zig hoặc Zag, khi nút cần đưa lên là con trực tiếp của gốc.
  \item Zig--Zig hoặc Zag--Zag, khi nút và cha cùng nằm về một phía.
  \item Zig--Zag hoặc Zag--Zig, khi nút và cha nằm khác phía.
\end{itemize}
Khác với cách quay từng bước đơn giản, splay ghép các phép quay theo cặp để
đảm bảo phân tích khấu hao tốt.

\section{Phân tích thế năng}

Slide nhắc các ký hiệu \(S(x)\), \(\rho(x)\) và \(\log_2 |S|\). Trong cách
phân tích chuẩn, \(S(x)\) là kích thước cây con của \(x\), còn hạng
\(\rho(x)=\lfloor \log_2 S(x)\rfloor\). Mỗi phép splay được trả bằng chi phí
thực cộng với thay đổi thế năng. Các bất đẳng thức giữa hạng của \(x,y,z\)
cho thấy tổng chi phí khấu hao của một lần splay bị chặn bởi hằng số nhân
\(\log n\).

\section{Ứng dụng}

Sau khi có thao tác splay, các phép tìm kiếm, chèn, xóa, tách và ghép cây đều
có thể xây dựng quanh ý tưởng ``đưa điểm nóng lên gốc''. Đây là lý do cây
splay phù hợp với các bài có chuỗi thao tác trực tuyến và mẫu truy cập thay
đổi theo thời gian.

\end{document}
